% sedphot User Manual
% Build with:  pdflatex sedphot_user_manual.tex   (run twice for the TOC)
%
% The format block below is shared verbatim with the sedfit manual.

\documentclass[11pt]{article}

\newcommand{\manualname}{sedphot}
\newcommand{\manualtagline}{Galaxy in, SED photometry out}
\newcommand{\manualversion}{0.3.0}

% ----------------------------------------------------------------
% FORMAT. This block is shared verbatim with the sibling user
% manual in the other repository so the two are typeset
% identically. Change it in both or in neither.
% ----------------------------------------------------------------
\usepackage[T1]{fontenc}
\usepackage{lmodern}
\usepackage[a4paper,margin=1in]{geometry}
\usepackage{fancyhdr}
\usepackage[dvipsnames]{xcolor}
\usepackage{booktabs}
\usepackage{longtable}
\usepackage{array}
\usepackage{enumitem}
\usepackage{titlesec}
\usepackage[most]{tcolorbox}
\usepackage[hidelinks]{hyperref}
\usepackage{url}

% --- Palette ---
\definecolor{slate}{RGB}{47,72,88}          % headings
\definecolor{boxrule}{RGB}{201,204,209}     % code-block border
\definecolor{boxfill}{RGB}{250,250,251}     % code-block fill
\definecolor{noteaccent}{RGB}{180,95,20}
\definecolor{warnaccent}{RGB}{160,30,30}

% --- Headings ---
\titleformat{\section}{\Large\bfseries\color{slate}}{\thesection}{1em}{}
\titleformat{\subsection}{\large\bfseries\color{slate}}{\thesubsection}{1em}{}

% --- Running head ---
\pagestyle{fancy}
\fancyhf{}
\lhead{\small\manualname\ user manual}
\rhead{\small\thepage}
\renewcommand{\headrulewidth}{0.4pt}

% --- Blocks ---
% A shell/code block: light fill, thin grey border, monospaced.
\newtcolorbox{cmd}{
  colback=boxfill, colframe=boxrule, boxrule=0.4pt,
  left=6pt, right=6pt, top=5pt, bottom=5pt,
  before skip=8pt, after skip=8pt, breakable
}

\newtcolorbox{notebox}[1][Note]{
  breakable, colback=noteaccent!5, colframe=noteaccent,
  boxrule=0.8pt, title=#1, fonttitle=\bfseries\small
}

\newtcolorbox{warnbox}[1][Careful]{
  breakable, colback=warnaccent!5, colframe=warnaccent,
  boxrule=0.8pt, title=#1, fonttitle=\bfseries\small
}

% One-line command.
\newcommand{\cmdline}[1]{\begin{cmd}\ttfamily #1\end{cmd}}

\setlist[itemize]{itemsep=2pt, topsep=4pt}
\setlist[enumerate]{itemsep=2pt, topsep=4pt}

% Table helpers: booktabs rules, wrapped description columns.
\newcolumntype{L}[1]{>{\raggedright\arraybackslash}p{#1}}
\newcolumntype{T}[1]{>{\ttfamily\raggedright\arraybackslash}p{#1}}

% --- Title page ---
\newcommand{\manualtitlepage}{%
  \thispagestyle{empty}
  \vspace*{4cm}
  \begin{center}
    {\Huge\bfseries\manualname}\\[0.7cm]
    {\Large User Manual}\\[1.3cm]
    {\normalsize\manualtagline}
  \end{center}
  \vfill
  \begin{center}\normalsize Version \manualversion\end{center}
  \newpage
  \tableofcontents
  \newpage
}

\begin{document}

\manualtitlepage

% ==========================================================
\section{What this does}
% ==========================================================

\texttt{sedphot} turns a galaxy into an SED-ready photometry table.

You give it a galaxy name (or a position on the sky) and an output
directory. It then does some or all of the following, depending on which
command you run:

\begin{itemize}\sloppy
  \item Looks the galaxy up in the public catalog archives (GALEX, SDSS,
        J-PLUS, Pan-STARRS, Legacy Surveys, AllWISE, HST) and records the
        published flux of the nearest source in every band.
  \item Downloads survey images (SDSS, Pan-STARRS, Legacy Surveys, CFHT,
        HST) and measures the galaxy itself in every band through one
        identical procedure, so the colors are internally consistent.
  \item Optionally requests SPHEREx spectrophotometry from IRSA.
  \item Writes a combined SED figure, quality-assurance figures, and a
        JSON record of exactly how every number was produced.
\end{itemize}

Fluxes come out in microjansky on the AB system. Everything for one
galaxy lands under one directory that you choose.

% ==========================================================
\section{Installing}
% ==========================================================

These instructions assume the machine has nothing set up yet.

\subsection{Get Python}

\texttt{sedphot} needs \textbf{Python 3.11 or newer} and a stack of
scientific libraries (NumPy, SciPy, Astropy, Astroquery, Matplotlib, and
others).

Use \textbf{Miniconda}. Conda installs Python and the scientific
libraries as pre-built binaries, which avoids the compiler toolchain that
some of these packages otherwise require, and it keeps each project's
libraries in a separate environment so one project cannot break another.

\textbf{Windows}
\begin{enumerate}\sloppy
  \item Go to \url{https://www.anaconda.com/download/success} (or
        \url{https://docs.conda.io/projects/miniconda/}) and download the
        \textbf{Miniconda Windows 64-bit installer} (\texttt{.exe}).
  \item Run the installer. Accept the defaults.
  \item From the Start menu, open \textbf{Anaconda Prompt}. Every command
        in this manual is typed there.
\end{enumerate}

\textbf{macOS}
\begin{enumerate}
  \item Download the \textbf{Miniconda macOS installer} for your chip
        (Apple silicon = \texttt{arm64}, Intel = \texttt{x86\_64}) from
        the same page.
  \item Run the installer, then open \textbf{Terminal}.
\end{enumerate}

\textbf{Linux}
\begin{enumerate}
  \item Download the Miniconda Linux \texttt{x86\_64} installer script
        from the same page.
  \item Run it with \texttt{bash Miniconda3-latest-Linux-x86\_64.sh},
        then open a new terminal.
\end{enumerate}

Check that conda is available:

\cmdline{conda -{}-version}

You should see something like \texttt{conda 24.9.2}.

\subsection{Get the code}

The source lives at \url{https://github.com/chrishopp12/query-photometry}.

\textbf{With git.} On Windows, install Git for Windows from
\url{https://git-scm.com/download/win} first; macOS and Linux usually
have git already. Then:

\cmdline{git clone https://github.com/chrishopp12/query-photometry.git}

This creates a folder named \texttt{query-photometry} in the current
directory. Move into it:

\cmdline{cd query-photometry}

\textbf{Without git.} Open
\url{https://github.com/chrishopp12/query-photometry} in a browser, click
the green \textbf{Code} button, choose \textbf{Download ZIP}, and unzip
it. You get a folder named \texttt{query-photometry-main}. Move into it:

\cmdline{cd query-photometry-main}

You are in the right folder if it contains a file named
\texttt{pyproject.toml}.

\subsection{Create an environment and install}

\textbf{Conda route (recommended, all platforms).}

\cmdline{conda create -n sedphot python=3.12}

Answer \texttt{y} when it asks. Then activate it:

\cmdline{conda activate sedphot}

Your prompt now starts with \texttt{(sedphot)}. Install the package:

\cmdline{pip install -e .}

That reads \texttt{pyproject.toml} and pulls in every dependency.

If any dependency fails to build (rare, but \texttt{reproject} is the
usual suspect), install the scientific stack from conda-forge first and
then install \texttt{sedphot} without re-resolving them:

\cmdline{conda install -c conda-forge numpy scipy pandas astropy}

\cmdline{conda install -c conda-forge astroquery matplotlib pillow}

\cmdline{conda install -c conda-forge reproject requests defusedxml}

\cmdline{\frenchspacing pip install -e . -{}-no-deps}

\textbf{pip / venv route (no conda).} Only use this if Python 3.11+ is
already installed and on your PATH.

Windows (Anaconda Prompt, PowerShell, or Command Prompt):

\cmdline{python -m venv .venv}

\cmdline{.venv\textbackslash Scripts\textbackslash activate}

macOS / Linux:

\cmdline{python3 -m venv .venv}

\cmdline{source .venv/bin/activate}

Then, on any platform:

\cmdline{pip install -e .}

\subsection{What ``editable'' means}

\texttt{pip install -e .} installs the package \textbf{in place}: Python
runs the source files in the folder you cloned, rather than a copy hidden
away in the environment. If you \texttt{git pull} a newer version, the
change takes effect immediately with no reinstall.

\subsection{Verify the install}

Three checks, in increasing strength.

\textbf{Is the command there?}

\cmdline{sedphot -{}-help}

Correct output begins:

\begin{cmd}\ttfamily\frenchspacing
usage: sedphot [-h]\\
\mbox{}~~~~~~~~~~~~~~~\{resolve,catalogs,measure,spherex,sed,overlay,remeasure,run\}\\
\mbox{}~~~~~~~~~~~~~~~...\\
\mbox{}\\
Galaxy in, SED photometry out: multi-archive retrieval and measurement.
\end{cmd}

\textbf{Is the right version installed?}

\cmdline{\frenchspacing python -c "import sedphot; print(sedphot.\_\_version\_\_)"}

Correct output:

\cmdline{0.3.0}

\textbf{Can it reach the network?} This resolves a real galaxy name
against the CDS Sesame service:

\cmdline{sedphot resolve -{}-name "NGC 4889"}

Correct output:

\begin{cmd}\ttfamily\frenchspacing
\mbox{}~~[resolve] 'NGC 4889' -> RA=195.033738, Dec=+27.977025~~(Sesame)\\
RA~~= 195.03373750\\
Dec = +27.97702500\\
label = NGC\_4889
\end{cmd}

If all three work, the install is good.

Every later session needs the environment activated first:

\cmdline{conda activate sedphot}

% ==========================================================
\section{Before you start}
% ==========================================================

The mental model is small.

\textbf{You pick a galaxy.} Either by name (\texttt{-{}-name "NGC 4889"},
resolved through Sesame, then NED, then SIMBAD) or by explicit position
(\texttt{-{}-ra 195.033738 -{}-dec 27.977025}). Give one or the other,
never both.

\textbf{You pick an output directory.} \texttt{-{}-out-dir} is required on
every command that writes anything. Everything for that galaxy lands under
it, inside a subdirectory called \texttt{Photometry/}: the tables, the
figures, the downloaded images, and the JSON records. One galaxy, one
directory.

\textbf{Files are named after a label.} The label defaults to the
sanitized name (\texttt{NGC 4889} becomes \texttt{NGC\_4889}) or, for a
bare position, an IAU-style designation
(\texttt{J130008.10+275837.3}). Override it with \texttt{-{}-label}. Output
files are \texttt{<label>\_catalog.csv},
\texttt{<label>\_measured.csv}, \texttt{<label>\_sed.png}, and so on.

\textbf{Downloads are cached.} The first run of a galaxy downloads images
and catalog queries; later runs on the same directory reuse them.
Re-running a command is cheap.

That is the whole model. A typical directory choice looks like:

\cmdline{sedphot run -{}-name "NGC 4889" -{}-out-dir NGC\_4889}

or, if you keep galaxies together:

Windows

\cmdline{sedphot run -{}-name "NGC 4889" -{}-out-dir galaxies\textbackslash NGC\_4889}

macOS / Linux

\cmdline{sedphot run -{}-name "NGC 4889" -{}-out-dir galaxies/NGC\_4889}

Path separators are the only difference between platforms in this manual.
Everything else is identical.

% ==========================================================
\section{The verbs}
% ==========================================================

\texttt{sedphot} has eight commands.

\begin{longtable}{T{2.2cm} L{4.6cm} L{7.2cm}}
\toprule
\textbf{\normalfont Verb} & \textbf{What it does} & \textbf{What it writes} \\
\midrule
\endhead
resolve & Turns a name into coordinates and an output label. &
  Nothing. Prints only. \\[2pt]
catalogs & Asks each catalog archive for its published photometry of the
  nearest source. & \texttt{<label>\_catalog.csv} + sidecar,
  \texttt{coverage\_catalogs.json} \\[2pt]
measure & Downloads survey images and measures the galaxy in every band.
  & \texttt{<label>\_measured.csv} + sidecar,
  \texttt{coverage\_measure.json}, cached images, per-band QA figures,
  \texttt{QA/growth\_curves.png} \\[2pt]
spherex & Requests a SPHEREx spectrophotometry table from IRSA. &
  \texttt{SPHEREx/table\_photometry.<tag>.csv} + sidecar,
  \texttt{SPHEREx/extractions.json} \\[2pt]
sed & Draws the combined SED from tables already on disk. &
  \texttt{<label>\_sed.png} \\[2pt]
overlay & Draws every archive's matched position on the HST color image.
  & \texttt{<label>\_overlay.png}, cached HST composite \\[2pt]
remeasure & Re-reports fluxes from a finished measurement at a different
  aperture. & Nothing, unless you pass \texttt{-{}-out} (a CSV) or
  \texttt{-{}-write-qa} (figures). \\[2pt]
run & \texttt{catalogs}, then \texttt{measure}, then optionally
  \texttt{spherex}, then \texttt{sed}. & Everything the stages above
  write. \\
\bottomrule
\end{longtable}

\texttt{-{}-out-dir} is \textbf{required} on \texttt{catalogs},
\texttt{measure}, \texttt{spherex}, \texttt{sed}, \texttt{overlay}, and
\texttt{run}. \texttt{resolve} takes no output directory because it writes
nothing. \texttt{remeasure} takes a path to a finished measurement's JSON
record instead.

% ==========================================================
\section{Your first run}
% ==========================================================

This walks one galaxy from nothing to a measured table, using only the
Legacy Surveys images so it finishes in a couple of minutes.

\textbf{Step 1.} Activate the environment.

\cmdline{conda activate sedphot}

\textbf{Step 2.} Confirm the galaxy resolves.

\cmdline{sedphot resolve -{}-name "NGC 4889"}

\begin{cmd}\ttfamily\frenchspacing
\mbox{}~~[resolve] 'NGC 4889' -> RA=195.033738, Dec=+27.977025~~(Sesame)\\
RA~~= 195.03373750\\
Dec = +27.97702500\\
label = NGC\_4889
\end{cmd}

\textbf{Step 3.} Fetch the catalog photometry.

\cmdline{sedphot catalogs -{}-name "NGC 4889" -{}-all -{}-out-dir NGC\_4889}

The console shows one block per archive, then a summary:

\begin{cmd}\ttfamily\frenchspacing
\mbox{}~~[resolve] 'NGC 4889' -> RA=195.033738, Dec=+27.977025~~(Sesame)\\
\mbox{}\\
Target: RA=195.033738, Dec=+27.977025~~~(search radius=2.0")\\
\mbox{}\\
=== galex ===\\
\mbox{}~~[GALEX GUVcat] Attempt 1/5, radius=2.0"\\
\mbox{}~~[GALEX GUVcat] Found 1 match(es) at radius=2.0"\\
\mbox{}\\
=== sdss ===\\
\mbox{}~~[SDSS] Attempt 1/5, radius=2.0"\\
\mbox{}~~[SDSS] Found 1 match(es) at radius=2.0"\\
\mbox{}\\
...\\
\mbox{}\\
Provider summary:\\
\mbox{}~~galex~~~~~~~~ok~~~~~~~~~~~~~~2 rows\\
\mbox{}~~sdss~~~~~~~~~ok~~~~~~~~~~~~~~5 rows\\
\mbox{}~~jplus~~~~~~~~no\_match~~~~~~~~0 rows -{}- no J-PLUS DR3 source found\\
\mbox{}~~panstarrs~~~~ok~~~~~~~~~~~~~~5 rows\\
\mbox{}~~legacy~~~~~~~ok~~~~~~~~~~~~~~7 rows\\
\mbox{}~~allwise~~~~~~ok~~~~~~~~~~~~~~4 rows\\
\mbox{}~~hst~~~~~~~~~~ok~~~~~~~~~~~~~~2 rows\\
\mbox{}\\
Saved 25 photometric points to: NGC\_4889/Photometry/NGC\_4889\_catalog.csv
\end{cmd}

followed by the table itself.

\textbf{Step 4.} Measure the images.

\cmdline{sedphot measure -{}-name "NGC 4889" -{}-instruments legacy -{}-out-dir NGC\_4889}

\begin{cmd}\ttfamily\frenchspacing
\mbox{}~~[resolve] 'NGC 4889' -> RA=195.033738, Dec=+27.977025~~(Sesame)\\
\mbox{}\\
Target: RA=195.033738, Dec=+27.977025~~~(aperture 12", scene engine)\\
\mbox{}\\
=== legacy images ===\\
\mbox{}~~3 band image(s) ready\\
\mbox{}\\
\mbox{}~~(scene catalog and star queries report here on a first run)\\
\mbox{}\\
\mbox{}~~~~[Legacy\_r] seat solve [target:sersic]: nfev 168, 4.1s last, ...\\
\mbox{}~~~~[Legacy\_r] background track: -0.0074 -> -0.0071 uJy/as2 (6/121 ...\\
\mbox{}~~~~[Legacy\_r] fill: 4\% of aperture masked, 100\% twin-filled, ...\\
\mbox{}~~~~[Legacy\_r] f\_ap = 30184.2 +- 41.6 uJy (empap), excess -22.4, ...\\
\mbox{}~~Legacy r: 30184.2 +/- 41.6 uJy (empap; QA Legacy\_r.png)\\
\mbox{}~~~~[Legacy\_z] ...\\
\mbox{}~~Legacy z: 48120.7 +/- 63.9 uJy (empap; QA Legacy\_z.png)\\
\mbox{}~~~~[Legacy\_g] ...\\
\mbox{}~~Legacy g: 12507.5 +/- 27.3 uJy (empap; QA Legacy\_g.png)\\
\mbox{}\\
Provider summary:\\
\mbox{}~~legacy~~~~~~~ok~~~~~~~~~~~~~~3 rows -{}- measured bands: r, z, g\\
\mbox{}\\
Saved 3 measured bands to: NGC\_4889/Photometry/NGC\_4889\_measured.csv
\end{cmd}

followed by the table. The numbers depend on the galaxy; the structure
does not. The indented \texttt{[Legacy\_r]} lines are the engine's running
commentary and are truncated above to fit the page; on screen they run to
the end of the line.

\textbf{Step 5.} Draw the SED.

\cmdline{sedphot sed -{}-out-dir NGC\_4889}

\cmdline{\frenchspacing \mbox{}~~[sed] wrote NGC\_4889/Photometry/NGC\_4889\_sed.png}

\textbf{The resulting file tree.}

\begin{cmd}\ttfamily\frenchspacing
NGC\_4889/\\
+-{}- Photometry/\\
\mbox{}~~~~+-{}- NGC\_4889\_catalog.csv\\
\mbox{}~~~~+-{}- NGC\_4889\_catalog.provenance.json\\
\mbox{}~~~~+-{}- NGC\_4889\_measured.csv\\
\mbox{}~~~~+-{}- NGC\_4889\_measured.provenance.json\\
\mbox{}~~~~+-{}- NGC\_4889\_sed.png\\
\mbox{}~~~~+-{}- coverage\_catalogs.json\\
\mbox{}~~~~+-{}- coverage\_measure.json\\
\mbox{}~~~~+-{}- Legacy/\\
\mbox{}~~~~|~~~+-{}- legacy\_ls-dr9\_g.fits\\
\mbox{}~~~~|~~~+-{}- legacy\_ls-dr9\_r.fits\\
\mbox{}~~~~|~~~+-{}- legacy\_ls-dr9\_z.fits\\
\mbox{}~~~~|~~~+-{}- QA/\\
\mbox{}~~~~|~~~~~~~+-{}- Legacy\_g.png\\
\mbox{}~~~~|~~~~~~~+-{}- Legacy\_r.png\\
\mbox{}~~~~|~~~~~~~+-{}- Legacy\_z.png\\
\mbox{}~~~~+-{}- QA/\\
\mbox{}~~~~|~~~+-{}- growth\_curves.png\\
\mbox{}~~~~+-{}- scene/\\
\mbox{}~~~~~~~~+-{}- gaia\_scene.csv\\
\mbox{}~~~~~~~~+-{}- tractor\_scene\_dr9.csv
\end{cmd}

Steps 3, 4 and 5 together are exactly what the \texttt{run} verb does:

\cmdline{sedphot run -{}-name "NGC 4889" -{}-out-dir NGC\_4889}

% ==========================================================
\section{Common tasks}
% ==========================================================

Each entry gives when to use it and the exact command. Substitute your own
galaxy name and output directory.

\subsection{A note on long commands}

A few commands below are too wide for one line and are shown broken across
two, with a trailing \texttt{\textbackslash}. That is the continuation
character for bash and zsh (macOS, Linux).

\begin{center}
\begin{tabular}{L{6.0cm} L{5.5cm}}
\toprule
\textbf{Shell} & \textbf{Continuation character} \\
\midrule
macOS / Linux (bash, zsh) & \texttt{\textbackslash} \\
Windows Command Prompt / Anaconda Prompt & \texttt{\textasciicircum} \\
Windows PowerShell & \texttt{\textasciigrave}\ (backtick) \\
\bottomrule
\end{tabular}
\end{center}

On Windows, either swap the trailing \texttt{\textbackslash} for the right
character, or --- simplest and always correct --- delete the line break
and paste the whole thing as one long line.

\subsection{Look up a galaxy's coordinates from its name}

Use this before anything else, to confirm the name resolves and to see
what label the output files will use.

\cmdline{sedphot resolve -{}-name "NGC 4889"}

If the name does not resolve, look the position up yourself (NED or SIMBAD
in a browser) and use \texttt{-{}-ra} / \texttt{-{}-dec} everywhere
instead:

\cmdline{sedphot resolve -{}-ra 195.033738 -{}-dec 27.977025}

\subsection{Fetch catalog photometry from every archive}

The normal starting point. Queries all seven catalog archives and writes
one combined table.

\cmdline{sedphot catalogs -{}-name "M87" -{}-all -{}-out-dir M87}

Each archive starts with a 2-arcsecond search radius and doubles it up to
five times if nothing is found. Raise the starting radius for a galaxy
whose catalog position is offset:

\cmdline{sedphot catalogs -{}-name "M87" -{}-all -{}-radius 5 -{}-out-dir M87}

\subsection{Fetch catalog photometry from one archive}

Use this when you only need one archive, or when one archive was down
during an earlier run and you want to retry just that one. The valid names
are \texttt{galex}, \texttt{sdss}, \texttt{jplus}, \texttt{panstarrs},
\texttt{legacy}, \texttt{allwise}, \texttt{hst}.

\cmdline{sedphot catalogs -{}-name "M87" -{}-instruments legacy -{}-out-dir M87}

Several at once:

\cmdline{sedphot catalogs -{}-name "M87" -{}-instruments legacy sdss -{}-out-dir M87}

Note that the table is rewritten from scratch each time, so a run with one
archive replaces a table that held all seven. Re-run with
\texttt{-{}-all} to rebuild the full table.

\subsection{Measure images from every survey}

Downloads images from all five image archives and measures the galaxy in
every band. This is the slow command: expect minutes per band on a first
run, less afterwards because the images are cached.

\cmdline{sedphot measure -{}-name "M87" -{}-all -{}-out-dir M87}

\subsection{Measure images from specific surveys only}

The valid names are \texttt{sdss}, \texttt{panstarrs}, \texttt{legacy},
\texttt{cfht}, \texttt{hst}.

\cmdline{sedphot measure -{}-name "M87" -{}-instruments legacy cfht -{}-out-dir M87}

\subsection{Measure a subset of filters}

\texttt{-{}-bands} applies the same filter list to every selected survey;
a survey that does not have one of them simply skips it.

\cmdline{sedphot measure -{}-name "M87" -{}-all -{}-bands g r z -{}-out-dir M87}

The default filter sets are:

\begin{center}
\begin{tabular}{L{5.2cm} L{7.6cm}}
\toprule
\textbf{Survey} & \textbf{Default filters} \\
\midrule
SDSS & \texttt{ugriz} \\
Pan-STARRS & \texttt{grizy} \\
Legacy Surveys (dr9) & \texttt{grz} \\
Legacy Surveys (dr10) & \texttt{griz} \\
CFHT MegaCam & \texttt{ugriz} \\
HST & every filter in the matched observation group \\
\bottomrule
\end{tabular}
\end{center}

For HST, filter names are flight-filter designations:

\begin{cmd}\ttfamily\frenchspacing
sedphot measure -{}-name "M87" -{}-instruments hst \textbackslash\\
\mbox{}~~~~-{}-bands F475W F814W -{}-out-dir M87
\end{cmd}

\subsection{Change the aperture size}

The measurement aperture is a circle centered on the target. The default
radius is \textbf{12 arcseconds}.

\cmdline{sedphot measure -{}-name "M87" -{}-all -{}-aperture 15 -{}-out-dir M87}

The aperture must fall inside the curve-of-growth radius grid, which by
default runs from 2 to 40 arcseconds. To go outside that, move the grid
too:

\begin{cmd}\ttfamily\frenchspacing
sedphot measure -{}-name "M87" -{}-all -{}-aperture 50 \textbackslash\\
\mbox{}~~~~-{}-radii 2 10 20 30 40 50 60 -{}-cutout-size 200 -{}-out-dir M87
\end{cmd}

\subsection{Re-report an existing measurement at a different aperture}

This is the fast path. Every measurement stores the \emph{curve of growth}
--- the enclosed flux at every radius from 2 to 40 arcseconds --- so
changing the reported aperture is a lookup, not a re-measurement. It takes
milliseconds and never touches the files on disk unless you ask it to.

Point it at the measurement's JSON record:

Windows

\begin{cmd}\ttfamily\frenchspacing
sedphot remeasure M87\textbackslash Photometry\textbackslash M87\_measured.provenance.json \textasciicircum\\
\mbox{}~~~~-{}-mode aperture -{}-aperture 10
\end{cmd}

macOS / Linux

\begin{cmd}\ttfamily\frenchspacing
sedphot remeasure M87/Photometry/M87\_measured.provenance.json \textbackslash\\
\mbox{}~~~~-{}-mode aperture -{}-aperture 10
\end{cmd}

The result prints as a table. To save it, add \texttt{-{}-out}:

\begin{cmd}\ttfamily\frenchspacing
sedphot remeasure M87/Photometry/M87\_measured.provenance.json \textbackslash\\
\mbox{}~~~~-{}-mode aperture -{}-aperture 10 -{}-out M87\_R10.csv
\end{cmd}

Two other useful forms. The total flux of the fitted model:

\begin{cmd}\ttfamily\frenchspacing
sedphot remeasure M87/Photometry/M87\_measured.provenance.json \textbackslash\\
\mbox{}~~~~-{}-integrated
\end{cmd}

And the fitted model's flux inside a chosen radius
(\texttt{-{}-mode sersic} is the default):

\begin{cmd}\ttfamily\frenchspacing
sedphot remeasure M87/Photometry/M87\_measured.provenance.json \textbackslash\\
\mbox{}~~~~-{}-aperture 20
\end{cmd}

Asking for an aperture \textbf{past} the stored grid in
\texttt{-{}-mode aperture} is different: there is no stored value to look
up, so the command rebuilds the scene from the saved fit and integrates.
It says so, and it takes minutes rather than milliseconds.

\subsection{Measure a small or compact target}

The defaults are sized for survey-resolution images of ordinary galaxies.
A compact target --- typically on an HST mosaic --- needs four things
moved together, or the measurement will treat the galaxy itself as sky.

\begin{center}
\begin{tabular}{T{3.2cm} L{3.6cm} L{4.6cm}}
\toprule
\textbf{\normalfont Flag} & \textbf{Default} & \textbf{Compact-target value} \\
\midrule
-{}-aperture & 12 arcsec & 2 arcsec \\
-{}-cutout-size & 120 arcsec & 20 arcsec \\
-{}-sky-rmin & 15 arcsec & 4 arcsec \\
-{}-radii & 2 to 40 arcsec & 0.5 to 8 arcsec \\
\bottomrule
\end{tabular}
\end{center}

\texttt{-{}-sky-rmin} is the boundary between ``target'' and ``sky''. Sky
is measured only outside it, and no pixel inside it is allowed to vote on
the background level. At the 15-arcsecond default, a compact HST target
would have most of the visit excluded from the sky estimate.

\begin{cmd}\ttfamily\frenchspacing
sedphot measure -{}-name "M87" -{}-instruments hst -{}-aperture 2 \textbackslash\\
\mbox{}~~~~-{}-cutout-size 20 -{}-sky-rmin 4 -{}-radii 0.5 1 1.5 2 3 4 6 8 \textbackslash\\
\mbox{}~~~~-{}-out-dir M87
\end{cmd}

The aperture must lie within the \texttt{-{}-radii} grid; the grid's outer
radius must be under half the cutout size.

\subsection{Run everything for one galaxy in a single command}

Catalogs, then images and measurement, then the SED plot. Every stage is
isolated: if one fails, the failure is recorded and the rest still run,
and the command exits with a nonzero status.

\cmdline{sedphot run -{}-name "NGC 4889" -{}-out-dir NGC\_4889}

Add SPHEREx:

\cmdline{sedphot run -{}-name "NGC 4889" -{}-spherex sersic -{}-out-dir NGC\_4889}

Leave out slow or irrelevant archives (\texttt{-{}-skip} accepts both
catalog and image archive names):

\cmdline{sedphot run -{}-name "NGC 4889" -{}-skip hst jplus -{}-out-dir NGC\_4889}

\begin{sloppypar}
\texttt{run} accepts every measurement flag \texttt{measure} does ---
\texttt{-{}-aperture}, \texttt{-{}-bands}, \texttt{-{}-cutout-size},
\texttt{-{}-sky-rmin}, \texttt{-{}-radii}, and the rest --- with one
exception: \texttt{-{}-dump-arrays} is available only on
\texttt{measure}.
\end{sloppypar}

\subsection{Make the SED plot from tables already on disk}

Use this after editing a table by hand, or after re-running one stage. It
reads whichever of \texttt{<label>\_catalog.csv} and
\texttt{<label>\_measured.csv} exist and draws them on one axis.

\cmdline{sedphot sed -{}-out-dir NGC\_4889}

If the directory holds tables for more than one label, name the one you
want:

\cmdline{sedphot sed -{}-label NGC\_4889 -{}-out-dir NGC\_4889}

\subsection{Check that the archives all matched the same object}

This is the check the flux columns cannot make: two archives confidently
reporting photometry for two different objects. It draws each archive's
matched position on the HST color image of the field, so a mismatch is
visible.

\cmdline{sedphot overlay -{}-out-dir NGC\_4889}

It only works where HST has imaged the field. The first run downloads a
\textbf{\textasciitilde380 MB} file (the HST detection mosaic, needed only
for its coordinate grid) and caches it. If you already have a FITS file on
the same grid, point at it and skip the download:

\begin{cmd}\ttfamily\frenchspacing
sedphot overlay -{}-out-dir NGC\_4889 \textbackslash\\
\mbox{}~~~~-{}-wcs-from NGC\_4889/Photometry/HST/hst\_F814W\_sci.fits
\end{cmd}

The file must be on the same pixel grid as the color composite; the
dimensions are checked, and a mismatch refuses rather than misplacing
every marker.

Adjust the panel sizes (both are half-widths in arcseconds):

\cmdline{sedphot overlay -{}-zoom-size 3 -{}-context-size 30 -{}-out-dir NGC\_4889}

\subsection{Fetch SPHEREx spectrophotometry}

Submits a forced-photometry job to the IRSA SPHEREx tool and writes the
raw per-visit table. The job runs on IRSA's servers; the command polls
until it finishes (up to an hour by default).

The default source model is an elliptical S\'ersic profile whose shape is
looked up in the Legacy Surveys Tractor catalog:

\cmdline{sedphot spherex -{}-name "NGC 4889" -{}-out-dir NGC\_4889}

If the shape lookup fails or returns nothing usable, the command
\textbf{stops} rather than guessing. Three ways forward:

Give the shape yourself, as n, axis ratio (a/b, at least 1), position
angle in degrees east of north, and effective radius in arcseconds:

\begin{cmd}\ttfamily\frenchspacing
sedphot spherex -{}-name "NGC 4889" -{}-out-dir NGC\_4889 \textbackslash\\
\mbox{}~~~~-{}-sersic-params 4 1.3 65 12
\end{cmd}

Fit the shape on an image instead:

\cmdline{sedphot spherex -{}-name "NGC 4889" -{}-sersic-from z -{}-out-dir NGC\_4889}

Or treat the target as a point source (this biases an extended source, in
a wavelength-dependent way):

\cmdline{sedphot spherex -{}-name "NGC 4889" -{}-model psf -{}-out-dir NGC\_4889}

Each distinct configuration gets its own tagged output file, so a PSF run
and a S\'ersic run coexist. Nothing is ever overwritten.

\subsection{Process a list of galaxies one after another}

Put one galaxy name per line in a plain text file,
\texttt{galaxies.txt}:

\begin{cmd}\ttfamily\frenchspacing
NGC 4889\\
NGC 4874\\
M87
\end{cmd}

\begin{sloppypar}
\textbf{Windows --- Command Prompt or Anaconda Prompt batch file.} Save as
\texttt{run\_all.bat} next to \texttt{galaxies.txt}, then double-click it
or type \texttt{run\_all.bat}:
\end{sloppypar}

\begin{cmd}\ttfamily\frenchspacing
@echo off\\
for /f "usebackq delims=" \%\%G in ("galaxies.txt") do (\\
\mbox{}~~echo === \%\%G ===\\
\mbox{}~~sedphot run -{}-name "\%\%G" -{}-out-dir "galaxies\textbackslash\%\%G"\\
)
\end{cmd}

\textbf{Windows --- PowerShell.} Type this at a PowerShell prompt, or save
it as \texttt{run\_all.ps1}:

\begin{cmd}\ttfamily\frenchspacing
Get-Content galaxies.txt | ForEach-Object \{\\
\mbox{}~~Write-Host "=== \$\_ ==="\\
\mbox{}~~sedphot run -{}-name "\$\_" -{}-out-dir "galaxies\textbackslash\$(\$\_ -replace ' ','\_')"\\
\}
\end{cmd}

\textbf{macOS / Linux --- bash or zsh.} Save as \texttt{run\_all.sh}, make
it executable with \texttt{chmod +x run\_all.sh}, then run
\texttt{./run\_all.sh}:

\begin{cmd}\ttfamily\frenchspacing
\#!/usr/bin/env bash\\
while IFS= read -r name; do\\
\mbox{}~~[ -z "\$name" ] \&\& continue\\
\mbox{}~~echo "=== \$name ==="\\
\mbox{}~~sedphot run -{}-name "\$name" -{}-out-dir "galaxies/\$\{name// /\_\}"\\
done < galaxies.txt
\end{cmd}

\begin{warnbox}
Run these one galaxy at a time, not in parallel. If you are also using
\texttt{-{}-registry-update} (see section 11), running in parallel will
corrupt the shared file.
\end{warnbox}

% ==========================================================
\section{Every flag, explained}
% ==========================================================

\subsection{Shared by \texttt{resolve}, \texttt{catalogs},
\texttt{measure}, \texttt{spherex}, \texttt{run}}

\begin{longtable}{T{2.9cm} L{2.15cm} L{6.45cm} L{2.7cm}}
\toprule
\textbf{\normalfont Flag} & \textbf{Takes} & \textbf{What it does} & \textbf{Default} \\
\midrule
\endhead
-{}-name & text & Resolvable object name; tried against Sesame, then NED,
  then SIMBAD. Use instead of \texttt{-{}-ra}/\texttt{-{}-dec}. & none \\[2pt]
-{}-ra & number & Right ascension in decimal degrees. Requires
  \texttt{-{}-dec}. & none \\[2pt]
-{}-dec & number & Declination in decimal degrees. Requires
  \texttt{-{}-ra}. & none \\[2pt]
-{}-out-dir & path & Directory everything lands under. \textbf{Required}
  on every verb except \texttt{resolve} and \texttt{remeasure}. & none \\[2pt]
-{}-label & text & Output filename stem. & sanitized name, or an IAU
  designation for a bare position \\
\bottomrule
\end{longtable}

Give exactly one of \texttt{-{}-name} or the
\texttt{-{}-ra}/\texttt{-{}-dec} pair.

\subsection{\texttt{catalogs}}

\begin{longtable}{T{2.9cm} L{2.15cm} L{6.45cm} L{2.7cm}}
\toprule
\textbf{\normalfont Flag} & \textbf{Takes} & \textbf{What it does} & \textbf{Default} \\
\midrule
\endhead
-{}-instruments & one or more names & Archives to query: \texttt{allwise},
  \texttt{galex}, \texttt{hst}, \texttt{jplus}, \texttt{legacy},
  \texttt{panstarrs}, \texttt{sdss}. & none \\[2pt]
-{}-all & (switch) & Query every archive. Use this or
  \texttt{-{}-instruments}. & off \\[2pt]
-{}-radius & arcsec & Starting cone-search radius. Doubles up to five
  times if nothing is found. & 2.0 \\[2pt]
-{}-legacy-dr & \texttt{dr9} or \texttt{dr10} & Legacy Surveys data
  release. \texttt{dr10} adds i-band in the south. & \texttt{dr9} \\[2pt]
-{}-dered & (switch) & Correct fluxes for Milky Way extinction. Off means
  as-measured, with the correction factor recorded per row. & off \\
\bottomrule
\end{longtable}

\subsection{\texttt{measure}}

\begin{longtable}{T{2.9cm} L{2.15cm} L{6.45cm} L{2.7cm}}
\toprule
\textbf{\normalfont Flag} & \textbf{Takes} & \textbf{What it does} & \textbf{Default} \\
\midrule
\endhead
-{}-instruments & one or more names & Image archives: \texttt{cfht},
  \texttt{hst}, \texttt{legacy}, \texttt{panstarrs}, \texttt{sdss}. &
  none \\[2pt]
-{}-all & (switch) & Every image archive. Use this or
  \texttt{-{}-instruments}. & off \\[2pt]
-{}-mode & \texttt{aperture} or \texttt{sersic} & \texttt{aperture}
  reports the flux inside the circular aperture. \texttt{sersic} reports
  the fitted galaxy model's flux instead. & \texttt{aperture} \\[2pt]
-{}-bands & one or more filters & Filter subset applied to every archive.
  & each archive's own default set \\[2pt]
-{}-aperture & arcsec & Aperture radius. Must lie inside the
  \texttt{-{}-radii} grid. & \textbf{12.0} \\[2pt]
-{}-radii & one or more arcsec & Radii at which the curve of growth is
  stored. & 2 to 40 in 1-arcsec steps \\[2pt]
-{}-cutout-size & arcsec & Width of the image cutout. Half of this is the
  largest usable radius. & 120 \\[2pt]
-{}-sky-rmin & arcsec & Boundary between target and sky. Sky is measured
  only outside it. Lower it for compact targets. & 15 \\[2pt]
-{}-registry & path & JSON file of shared neighbor shapes to reuse
  (section 11). & none \\[2pt]
-{}-registry-\newline update & (switch) & Also write this galaxy's results
  back into \texttt{-{}-registry}. Requires \texttt{-{}-registry}. &
  off \\[2pt]
-{}-sersic-from & band & With \texttt{-{}-mode sersic}: fit the galaxy's
  shape on this band (\texttt{z} or \texttt{Legacy\_z}) and use it for all
  bands. & per-archive refit \\[2pt]
-{}-sersic-\newline params & 4 numbers & With \texttt{-{}-mode sersic}: state the
  shape explicitly as n, axis ratio (a/b >= 1), position angle (deg E of
  N), effective radius (arcsec). & none \\[2pt]
-{}-sersic-\newline seeing & arcsec & PSF full width at half maximum assumed by a
  \texttt{-{}-sersic-from} fit. The fitted index and radius are sensitive
  to it. & a typical value for that survey, not a measured one; the run
  prints a warning saying so \\[2pt]
-{}-legacy-dr & \texttt{dr9} or \texttt{dr10} & Legacy Surveys release for
  both images and the scene catalog. & \texttt{dr9} \\[2pt]
-{}-legacy-\newline bricks & (switch) & Download full Legacy brick coadds instead
  of small cutouts. Adds real per-pixel errors; roughly 40 MB per file. &
  off \\[2pt]
-{}-hst-\newline proposal-id & text & Restrict HST to one observing
  program. & none \\[2pt]
-{}-dump-arrays & (switch) & Write raw per-band arrays under
  \texttt{<Instrument>/QA/} for debugging. & off \\
\bottomrule
\end{longtable}

\begin{sloppypar}
\texttt{-{}-sersic-from}, \texttt{-{}-sersic-params} and
\texttt{-{}-sersic-seeing} are refused under
\texttt{-{}-mode aperture}; the command tells you so rather than ignoring
them.
\end{sloppypar}

\subsection{\texttt{spherex}}

\begin{longtable}{T{2.9cm} L{2.15cm} L{6.45cm} L{2.7cm}}
\toprule
\textbf{\normalfont Flag} & \textbf{Takes} & \textbf{What it does} & \textbf{Default} \\
\midrule
\endhead
-{}-model & \texttt{psf} or \texttt{sersic} & Source model for the forced
  photometry. \texttt{psf} biases extended sources. & \texttt{sersic} \\[2pt]
-{}-sersic-\newline params & 4 numbers & Explicit shape: n (<= 6), axis ratio
  (a/b >= 1), position angle (deg E of N), effective radius (arcsec). &
  Tractor catalog lookup \\[2pt]
-{}-sersic-from & band & Fit the shape on this band's image instead of the
  catalog lookup. & none \\[2pt]
-{}-sersic-\newline seeing & arcsec & PSF full width at half maximum of the
  shape-fit band. & a typical value for that survey, not a measured
  one \\[2pt]
-{}-bkg-size & pixels & Size of the tool's background-estimation region. &
  15 \\[2pt]
-{}-mjd-range & 2 numbers & Restrict to visits inside this MJD window.
  This is IRSA's documented workaround for epochs whose metadata breaks
  the job. & none \\[2pt]
-{}-poll & seconds & How often to ask IRSA whether the job is done. & 5 \\[2pt]
-{}-timeout & seconds & Give up after this long. & 3600 \\[2pt]
-{}-legacy-dr & \texttt{dr9} or \texttt{dr10} & Legacy release for the
  shape lookup or shape-fit image. & \texttt{dr9} \\
\bottomrule
\end{longtable}

\subsection{\texttt{sed}}

\begin{longtable}{T{2.9cm} L{2.15cm} L{6.45cm} L{2.7cm}}
\toprule
\textbf{\normalfont Flag} & \textbf{Takes} & \textbf{What it does} & \textbf{Default} \\
\midrule
\endhead
-{}-out-dir & path & Directory holding the tables. Required. & none \\[2pt]
-{}-label & text & Which table family to plot. & inferred when the
  directory holds exactly one \\
\bottomrule
\end{longtable}

\subsection{\texttt{overlay}}

\begin{longtable}{T{2.9cm} L{2.15cm} L{6.45cm} L{2.7cm}}
\toprule
\textbf{\normalfont Flag} & \textbf{Takes} & \textbf{What it does} & \textbf{Default} \\
\midrule
\endhead
-{}-out-dir & path & Directory holding the tables. Required. & none \\[2pt]
-{}-label & text & Which table family to plot. & inferred when
  unambiguous \\[2pt]
-{}-zoom-size & arcsec & Half-width of the zoom panel. & 5 \\[2pt]
-{}-context-size & arcsec & Half-width of the context panel. & 15 \\[2pt]
-{}-wcs-from & path & A local FITS file already on the composite's grid,
  used instead of downloading the \textasciitilde380 MB detection mosaic.
  Refuses if the path names no file, or if the grid does not match. &
  none \\[2pt]
-{}-dpi & number & Figure resolution. & 200 \\
\bottomrule
\end{longtable}

\subsection{\texttt{remeasure}}

The first argument is positional: the path to a
\texttt{<label>\_measured.provenance.json}.

\begin{longtable}{T{2.9cm} L{2.15cm} L{6.45cm} L{2.7cm}}
\toprule
\textbf{\normalfont Flag} & \textbf{Takes} & \textbf{What it does} & \textbf{Default} \\
\midrule
\endhead
-{}-mode & \texttt{sersic} or \texttt{aperture} & \texttt{sersic} reports
  the fitted model's flux; \texttt{aperture} reports the measured,
  neighbor-subtracted flux. & \texttt{sersic} \\[2pt]
-{}-aperture & arcsec & Radius to report at. Matches the
  \texttt{measure} default, so a bare \texttt{remeasure} re-reports at the
  radius the measurement used. & 12 \\[2pt]
-{}-integrated & (switch) & Report the total instead of an aperture. In
  \texttt{sersic} mode that is the model's total; in \texttt{aperture}
  mode it is the outermost stored radius. Overrides
  \texttt{-{}-aperture}. & off \\[2pt]
-{}-shape & \texttt{forced} or \texttt{fitted} & Which galaxy shape the
  report is built on. \texttt{forced} uses the archive's reference-band
  shape. \texttt{fitted} uses each band's own shape where one was stored;
  bands without one fall back to \texttt{forced} and say so. &
  \texttt{forced} \\[2pt]
-{}-registry & path & Registry file to fall back on when rebuilding past
  the stored radius grid. & a \texttt{registry.json} beside the galaxy
  directory, if present \\[2pt]
-{}-write-qa & (switch) & When rebuilding past the grid, write per-band
  figures to \texttt{QA/remeasure\_R<N>as/}. & off \\[2pt]
-{}-out & path & Write the result to this CSV instead of printing it. &
  print \\
\bottomrule
\end{longtable}

\subsection{\texttt{run}}

\texttt{run} takes every flag from \texttt{catalogs}
(\texttt{-{}-radius}, \texttt{-{}-dered}, \texttt{-{}-legacy-dr}) and
every flag from \texttt{measure} \textbf{except}
\texttt{-{}-instruments}, \texttt{-{}-all} and
\texttt{-{}-dump-arrays}, plus:

\begin{longtable}{T{2.9cm} L{2.15cm} L{6.45cm} L{2.7cm}}
\toprule
\textbf{\normalfont Flag} & \textbf{Takes} & \textbf{What it does} & \textbf{Default} \\
\midrule
\endhead
-{}-skip & one or more names & Archives to leave out. Accepts both catalog
  and image archive names. & none \\[2pt]
-{}-spherex & \texttt{off}, \texttt{psf}, \texttt{sersic} & Also fetch
  SPHEREx, with this source model. & off \\
\bottomrule
\end{longtable}

\texttt{run} uses every archive that is not skipped, which is why it has
no \texttt{-{}-instruments}.

One special case: \texttt{-{}-sersic-params} is normally refused under
\texttt{-{}-mode aperture}, but under \texttt{run} with
\texttt{-{}-spherex psf} or \texttt{-{}-spherex sersic} it is accepted,
because there it also declares the shape SPHEREx extracts with. One shape
per galaxy, shared by both stages.

% ==========================================================
\section{What comes out}
% ==========================================================

\subsection{The directory tree}

A full run of one galaxy produces:

\begin{cmd}\ttfamily\frenchspacing
<out-dir>/\\
+-{}- patches.json~~~~~~~~~~~~~~~~~~~~(optional, you create it)\\
+-{}- Photometry/\\
\mbox{}~~~~+-{}- <label>\_catalog.csv~~~~~~~~~~~~~~~~~~~~~archive photometry\\
\mbox{}~~~~+-{}- <label>\_catalog.provenance.json\\
\mbox{}~~~~+-{}- <label>\_measured.csv~~~~~~~~~~~~~~~~~~~~measured photometry\\
\mbox{}~~~~+-{}- <label>\_measured.provenance.json\\
\mbox{}~~~~+-{}- <label>\_sed.png~~~~~~~~~~~~~~~~~~~~~~~~~combined SED figure\\
\mbox{}~~~~+-{}- <label>\_overlay.png~~~~~~~~~~~~~~~~~~~~~position check figure\\
\mbox{}~~~~+-{}- coverage\_catalogs.json~~~~~~~~~~~~~~~~~~per-archive outcome\\
\mbox{}~~~~+-{}- coverage\_measure.json~~~~~~~~~~~~~~~~~~~per-archive outcome\\
\mbox{}~~~~+-{}- Legacy/~~~~~~~cached images + QA/<band>.png\\
\mbox{}~~~~+-{}- SDSS/~~~~~~~~~cached images + QA/<band>.png\\
\mbox{}~~~~+-{}- PanSTARRS/~~~~cached images + QA/<band>.png\\
\mbox{}~~~~+-{}- CFHT/~~~~~~~~~cached images + QA/<band>.png\\
\mbox{}~~~~+-{}- HST/~~~~~~~~~~cached mosaics + QA/<band>.png\\
\mbox{}~~~~+-{}- QA/\\
\mbox{}~~~~|~~~+-{}- growth\_curves.png~~~~~~~~~~~~~~~~~~~all bands, one figure\\
\mbox{}~~~~+-{}- scene/~~~~~~~~~~~~~~~~~~~~~~~~~~~~~~~~~~cached catalog queries\\
\mbox{}~~~~+-{}- SPHEREx/\\
\mbox{}~~~~~~~~+-{}- table\_photometry.<tag>.csv\\
\mbox{}~~~~~~~~+-{}- table\_photometry.<tag>.provenance.json\\
\mbox{}~~~~~~~~+-{}- extractions.json
\end{cmd}

\subsection{What each file is}

\begin{longtable}{T{5.4cm} L{8.8cm}}
\toprule
\textbf{\normalfont File} & \textbf{What it is} \\
\midrule
\endhead
<label>\_catalog.csv & One row per band, from the published archive
  catalogs. The nearest source to the requested position in each
  archive. \\[2pt]
<label>\_measured.csv & One row per band, measured by \texttt{sedphot}
  from the images. Same column layout as the catalog table. \\[2pt]
*.provenance.json & The complete record of how the neighboring file was
  made: every setting, the code revision, a checksum of the file, and for
  measurements the full fit. This is what \texttt{remeasure} reads. Do not
  edit it. \\[2pt]
<label>\_sed.png & Flux against wavelength, both tables together. Catalog
  points are filled, measured points are open. Marker shape identifies the
  instrument, color tracks wavelength. \\[2pt]
<label>\_overlay.png & Two panels of the HST color image with each
  archive's matched position marked. \\[2pt]
coverage\_catalogs.json,\newline coverage\_measure.json & Per-archive
  outcome: \texttt{ok}, \texttt{no\_coverage}, \texttt{no\_match}, or
  \texttt{error}, with a message and the number of rows. \\[2pt]
<Instrument>/*.fits & The downloaded images, cached. Deleting one forces a
  re-download on the next run. \\[2pt]
<Instrument>/QA/<band>.png & The per-band diagnostic figure
  (section 9). \\[2pt]
QA/growth\_curves.png & Enclosed flux against radius for every measured
  band on one axis. \\[2pt]
scene/*.csv & Cached catalog queries describing the neighboring sources.
  Deleting these forces a re-query. \\[2pt]
SPHEREx/table\_photometry.\newline <tag>.csv & The raw per-visit SPHEREx
  table, exactly as IRSA returned it. Never overwritten. \\[2pt]
SPHEREx/extractions.json & An index mapping each tag to the configuration
  that produced it. \\
\bottomrule
\end{longtable}

\subsection{The output table columns}

Both \texttt{<label>\_catalog.csv} and \texttt{<label>\_measured.csv}
share this layout.

\begin{longtable}{T{3.1cm} L{8.1cm} L{3.2cm}}
\toprule
\textbf{\normalfont Column} & \textbf{Meaning} & \textbf{Units} \\
\midrule
\endhead
band & Band label, \texttt{<Instrument>\_<filter>} --- e.g.\
  \texttt{Legacy\_g}, \texttt{SDSS\_u}, \texttt{PS1\_z},
  \texttt{CFHT\_r}, \texttt{GALEX\_FUV}, \texttt{WISE\_W1},
  \texttt{JPLUS\_uJAVA}, \texttt{HST\_F475W}. & --- \\[2pt]
flux\_uJy & The flux. & microjansky (AB) \\[2pt]
flux\_err\_uJy & Uncertainty on the flux. & microjansky \\[2pt]
mag\_AB & The same measurement as an AB magnitude. Empty when the flux is
  zero or negative. & mag \\[2pt]
mag\_err & Uncertainty on the magnitude. & mag \\[2pt]
target\_ra & The position you asked for. & degrees \\[2pt]
target\_dec & The position you asked for. & degrees \\[2pt]
match\_ra & Where the archive's matched source actually is. For measured
  rows this equals \texttt{target\_ra}. & degrees \\[2pt]
match\_dec & Where the archive's matched source actually is. & degrees \\[2pt]
sep\_arcsec & Distance from the requested position to the match. 0 for
  measured rows. & arcsec \\[2pt]
flags & Quality tokens, \texttt{key=value} separated by semicolons (see
  below). For catalog rows, whatever flags that archive publishes. &
  --- \\[2pt]
source & Where the number came from --- \texttt{Legacy\_DR9},
  \texttt{SDSS\_DR17\_cModel}, \texttt{PanSTARRS\_DR1},
  \texttt{GALEX\_GUVcat\_AIS}, \texttt{sedphot\_aperture\_scene\_empap},
  and so on. & --- \\[2pt]
retrieved & Date the value was pulled. & ISO date \\[2pt]
mw\_transmission & Milky Way transmission in this band, where the archive
  supplies it. Empty when unknown. & fraction \\[2pt]
dered\_applied & \texttt{True} if the flux has been corrected for Milky
  Way extinction. Measured rows are always \texttt{False}. & --- \\
\bottomrule
\end{longtable}

\textbf{Three things to know about the numbers.}

\begin{itemize}
  \item \textbf{Fluxes are microjansky on the AB system.} Convert to AB
        magnitude with \texttt{mag = 23.9 - 2.5 log10(flux)}.
  \item \textbf{The errors are statistical only.} They describe the noise
        in the measurement and nothing else --- no calibration
        uncertainty, no error floor, no model uncertainty. Adding those is
        the SED fitter's job, not this table's.
  \item \textbf{Negative fluxes are real.} A non-detection in a faint band
        gives a small negative number after background subtraction. That
        is the honest measurement and it is preserved deliberately; it is
        not an error, and it must not be clipped to zero. Its
        \texttt{mag\_AB} cell is simply empty.
\end{itemize}

\textbf{The \texttt{flags} tokens on measured rows.} These let you filter a
table without opening a single figure.

\begin{longtable}{T{2.6cm} L{5.4cm} L{6.4cm}}
\toprule
\textbf{\normalfont Token} & \textbf{Meaning} & \textbf{What to look for} \\
\midrule
\endhead
cov & Fraction of aperture pixels with real data. & Should be 1.000. Below
  0.95 the band is refused outright. \\[2pt]
maskfrac & Fraction of the aperture masked because a neighboring source
  sat there. & A few percent is normal. Above
  \textasciitilde0.2 the flux leans heavily on reconstruction. \\[2pt]
twinfrac & Fraction of masked pixels reconstructed from the galaxy's
  mirror side rather than from the model. & 1.00 is best. \\[2pt]
nbsub & Flux subtracted from the aperture as neighboring sources. &
  Compare against the reported flux. Large values mean a crowded
  field. \\[2pt]
excess & Flux growing outside the aperture beyond what the galaxy's own
  model predicts. & Near zero is good. Large positive means contamination
  or an unmodeled envelope. \\[2pt]
pedb & Residual uniform background left in the curve of growth. & Should
  be near zero. \\[2pt]
conv & Radius, in arcsec, at which the curve of growth flattened.
  \texttt{0} means it never did. & Should be at or inside the
  aperture. \\[2pt]
bg & The fitted background level. & --- \\[2pt]
far & An independent background level measured far from the target. &
  Should broadly agree with \texttt{bg}. \\[2pt]
mesh & Flux the smooth background surface removed from the aperture. &
  Large values mean structured background. \\[2pt]
art & Area, in square arcsec, flagged as a detector artifact (a bleed
  trail, a streak). & Present only when one was found. \\[2pt]
leash & How many neighboring sources hit their brightness limit. & Present
  only when some did. \\[2pt]
refit & The galaxy's fitted brightness relative to the catalog's. & Near 1
  is unremarkable. \\[2pt]
atbound & How many fitted shape parameters ended at a limit. & Present
  only when some did. \\[2pt]
reg & How many neighbors came from the shared registry. & --- \\[2pt]
scene=none & No catalog of neighboring sources was available; the galaxy
  was measured with no neighbor modeling at all. & Treat the flux with
  care in a crowded field. \\
\bottomrule
\end{longtable}

% ==========================================================
\section{Reading the QA figures}
% ==========================================================

Two figures per measurement run, plus one per band.

\subsection{The per-band figure}

\texttt{<Instrument>/QA/<Instrument>\_<band>.png} --- five panels in a row.

\begin{enumerate}
  \item \textbf{data} --- the image, with stars already subtracted. The
        blue circle is your aperture. Orange contours, if present, mark
        detector artifacts. The title gives the fitted background level
        and its tilt.
  \item \textbf{fitted scene + background} --- the model: the galaxy, its
        neighbors, and the background, as the fit believes them to be.
  \item \textbf{residual - mesh} --- the data minus the model, and minus
        the smooth background surface. This is the panel that tells you
        whether the fit worked.
  \item \textbf{masked + filled - mesh} --- the image the flux is actually
        measured from: masked pixels replaced by reconstructed values.
        Green contours outline the mask. The title gives the masked
        fraction of the aperture.
  \item \textbf{curve of growth} --- enclosed flux against radius. The red
        line with dots is the measurement; the black dashed line is the
        fitted galaxy model. The blue vertical line is your aperture; a
        grey dotted line marks where the curve flattened. The title gives
        the final flux, its error, and the excess growth.
\end{enumerate}

\textbf{A good figure.}

\begin{itemize}
  \item Panel 3 is flat noise. No bright leftover at the galaxy's center,
        no dark hole where a neighbor was over-subtracted.
  \item Panel 5's red curve rises and then flattens, and the flattening
        happens at or inside the blue aperture line. The title says
        \texttt{conv@14"} or similar rather than \texttt{no plateau}.
  \item The red curve and the black dashed model track each other out to
        the aperture.
  \item \texttt{excess} in the title is small compared with the flux.
  \item Green mask contours in panel 4 cover only a small part of the blue
        circle.
\end{itemize}

\textbf{A suspicious figure.}

\begin{longtable}{L{5.2cm} L{9.2cm}}
\toprule
\textbf{What you see} & \textbf{What it means} \\
\midrule
\endhead
Title says \texttt{no plateau} & The curve of growth was still rising at
  40 arcsec. Either the galaxy is genuinely large, or the background is
  under-subtracted. Check \texttt{pedb} and the \texttt{far}
  token. \\[2pt]
Red curve turns over and falls & The background was over-subtracted. \\[2pt]
Large positive \texttt{excess} & Light in the outskirts that the galaxy's
  model does not account for --- a neighbor, or a real envelope. \\[2pt]
Dark hole at a neighbor's position in panel 3 & That neighbor was
  over-subtracted. \\[2pt]
Bright residual core in panel 3 & The galaxy profile did not fit; the flux
  may be biased. \\[2pt]
Green contours across much of the blue circle & Most of the aperture is
  reconstructed rather than measured. Check \texttt{twinfrac}. \\[2pt]
Orange artifact contours crossing the aperture & A detector artifact sits
  on the target. \\[2pt]
Red and black curves diverge inside the aperture & The fitted model is a
  poor description of the galaxy. Fine in \texttt{-{}-mode aperture} (the
  model is only a cross-check there), serious in
  \texttt{-{}-mode sersic} (the model \emph{is} the answer). \\
\bottomrule
\end{longtable}

\subsection{The growth-curve figure}

\texttt{QA/growth\_curves.png} --- every measured band on one axis,
enclosed flux against radius, with a logarithmic vertical scale. Line style
identifies the instrument, color tracks wavelength. The blue vertical line
is the science aperture.

Read it for consistency across bands. Curves should be roughly parallel
and should flatten at similar radii. A single band that keeps climbing
while its neighbors flatten, or one that crosses the others, is the band
to open individually.

% ==========================================================
\section{When something goes wrong}
% ==========================================================

Every command reports failures rather than hiding them. Most messages name
the fix. The recurring ones:

\begin{longtable}{T{3.5cm} L{4.2cm} L{6.9cm}}
\toprule
\textbf{\normalfont Symptom} & \textbf{Cause} & \textbf{Fix} \\
\midrule
\endhead
error: the following arguments are required: -{}-out-dir &
  \texttt{-{}-out-dir} is required on every verb that writes. & Add
  \texttt{-{}-out-dir <directory>}. \\[2pt]
error: give -{}-instruments ... or -{}-all & \texttt{catalogs} and
  \texttt{measure} need to be told which archives to use. & Add
  \texttt{-{}-all}, or \texttt{-{}-instruments legacy sdss}. \\[2pt]
Could not resolve 'X' with Sesame, NED, or SIMBAD & The name is not in any
  resolver. & Look the position up yourself and use \texttt{-{}-ra} /
  \texttt{-{}-dec}. \\[2pt]
-{}-registry-update needs -{}-registry PATH & You asked to write to a
  registry without naming one. & Add \texttt{-{}-registry <path>}, or drop
  \texttt{-{}-registry-update}. \\[2pt]
-{}-sersic-from only applies to -{}-mode sersic & A shape flag was given
  under aperture mode, where it would have no effect. & Add
  \texttt{-{}-mode sersic}, or drop the flag. \\[2pt]
ValueError: -{}-aperture 60" lies outside the curve-of-growth grid
  [2, 40]" & The aperture is outside the radius grid the run stores. &
  Pass \texttt{-{}-radii} covering it, and a \texttt{-{}-cutout-size} at
  least twice the largest radius. \\[2pt]
{\normalfont An \textbf{image} archive reports} no\_coverage & That
  survey's footprint does not include this position. Nothing is wrong. &
  Nothing to do. Fewer bands is the correct answer. \\[2pt]
{\normalfont A \textbf{catalog} archive reports} no\_match & Either the
  field genuinely holds no catalogued source within the expanded search
  radius (including ``this position is outside that survey's footprint''),
  \textbf{or} the archive was down. The two are indistinguishable: a
  failed query returns no rows, and the radius expansion then reports
  \texttt{no\_match}. Catalog archives never report
  \texttt{no\_coverage}. & Re-run just that archive later. If it still
  says \texttt{no\_match}, try a larger \texttt{-{}-radius}. An unexpected
  \texttt{no\_match} on a well-covered galaxy is almost always an
  outage. \\[2pt]
{\normalfont An archive reports} error & A service or parsing failure that
  got past the provider's own retries. The message names it. & Read
  \texttt{coverage\_catalogs.json} or \texttt{coverage\_measure.json},
  then re-run that archive. \\[2pt]
Legacy z: no\_coverage -{}- aperture coverage 0.83 < 0.95 & More than 5\%
  of the aperture is off the image footprint, blank, or covered by an
  artifact. The band is refused rather than reporting a silently biased
  flux. & Use a smaller \texttt{-{}-aperture}, or a larger
  \texttt{-{}-cutout-size} if the cutout was simply too small. Otherwise
  accept the missing band. \\[2pt]
blank pixels inside the 0.5" peak & Dead pixels at the very center of the
  galaxy. No reconstruction can invent the peak. & The band cannot be
  measured. \\[2pt]
scene catalog is empty here; measuring blind, and scene=none in flags & No
  catalog of neighboring sources covers this position, so nothing was
  subtracted or masked. & The flux is still measured, but with no neighbor
  handling. In an uncrowded field this is fine; in a crowded one, treat it
  with care. \\[2pt]
No images fetched from any provider; nothing to measure. & No archive
  covered the position, or all of them failed. & Check
  \texttt{coverage\_measure.json} for which and why. \\[2pt]
{\normalfont \texttt{overlay} downloads \textasciitilde380 MB} & The HST
  color composite carries no coordinate grid of its own, so the detection
  mosaic is fetched for its grid. It is cached after the first time. &
  Pass \texttt{-{}-wcs-from} with a local FITS file already on that
  grid. \\[2pt]
no HAP total/detection product at this position, {\normalfont or}
  sedphot overlay: no figure written & HST has not imaged this field. &
  The overlay is unavailable for this galaxy. \\[2pt]
-{}-wcs-from names no file & The path is mistyped. It refuses rather than
  falling back to the big download. & Correct the path. \\[2pt]
[spherex] Tractor shape lookup failed ... aborting {\normalfont or}
  no usable extended Tractor shape at this position & SPHEREx will not
  guess a source shape. A raw table written under the wrong shape is
  permanent and unmarked, so it stops instead. & Pass
  \texttt{-{}-sersic-params N AXR PA RE}, or \texttt{-{}-sersic-from z} to
  fit one, or \texttt{-{}-model psf}. \\[2pt]
table\_photometry.\newline <tag>.csv exists but no sidecar vouches for
  this configuration & A raw SPHEREx table is on disk that nothing
  certifies. Nothing is ever overwritten. & Move the file aside
  deliberately, then re-run. \\[2pt]
[spherex] WARNING: reff 0.60" is below the tool's 1" point-source
  threshold & The tool treats anything that small as a point source, so
  the S\'ersic parameters do nothing. & Either accept it, or use
  \texttt{-{}-model psf} and be explicit. \\[2pt]
remeasure prints aperture 45" is past the stored curve of growth (ends at
  40"); rebuilding the scene & Beyond the stored grid there is no value to
  interpolate, so the scene is rebuilt and re-integrated. & Nothing wrong
  --- just expect minutes rather than milliseconds. \\[2pt]
sedphot remeasure: aperture 70" exceeds the stamp half-width (60") & The
  requested radius runs off the cutout the fit was built on. & Re-measure
  with a larger \texttt{-{}-cutout-size}. \\[2pt]
remeasure says bands fell back to the forced shape under -{}-shape fitted
  & Those bands have no per-band shape stored --- either an older
  measurement, or a galaxy where the extra fit never ran. & Either accept
  \texttt{forced}, or re-measure the galaxy. \\[2pt]
cannot infer -{}-label in ...: found stems [...] & The directory holds
  tables for more than one galaxy, so \texttt{sed} and \texttt{overlay}
  cannot guess which. & Pass \texttt{-{}-label <stem>}. \\[2pt]
[sed] WARNING ... the plotted fluxes are NOT on a common extinction scale
  & You used \texttt{-{}-dered}, which corrects the catalog table only.
  The measured table is always as-measured, so the figure mixes the two. &
  Expected. Do not read colors off that figure; use the tables. \\
\bottomrule
\end{longtable}

% ==========================================================
\section{Getting more out of it}
% ==========================================================

Two features exist for cases the defaults do not cover. You will not need
either at first.

\subsection{\texttt{patches.json} --- correcting the scene by hand}

Before measuring, \texttt{sedphot} builds a picture of the field from the
Legacy Surveys catalog: which sources are there, how big and how bright.
Occasionally that catalog is wrong in a way that matters --- most often it
merges a close pair into a single entry, or splits one galaxy into
several.

You can correct it. Put a file named \texttt{patches.json} in the galaxy's
\texttt{-{}-out-dir} (next to \texttt{Photometry/}, not inside it). It is
read automatically and its keys are logged at the start of the run.

The keys it accepts:

\begin{longtable}{T{3.6cm} L{10.4cm}}
\toprule
\textbf{\normalfont Key} & \textbf{What it does} \\
\midrule
\endhead
replace\_rows & Replace one catalog entry with one or more of your own.
  The usual fix for a merged close pair. \\[2pt]
free\_seats & Let the source at a given position have its shape fitted
  rather than taken from the catalog. \\[2pt]
target\_system & Declare that a nearby source belongs to the target ---
  its light is counted with the galaxy rather than subtracted from
  it. \\[2pt]
target\_halo & Give the target an extended-envelope component, for a
  brightest-cluster galaxy. \\[2pt]
target\_refit & Turn the target's own shape refit on or off. \\[2pt]
snap\_gated & Move fitted sources onto their nearest brightness peak. \\[2pt]
harvest\_target & Record the target's own shape into the shared
  registry. \\[2pt]
comment & Free text. Ignored by the code; use it to say why the patch
  exists. \\
\bottomrule
\end{longtable}

Reach for this when a QA figure shows a clear over- or under-subtraction
at a specific neighbor, and only then. An unnecessary patch is a way to
bias a measurement.

\subsection{The cross-field registry --- shared neighbors}

When several galaxies in the same field are measured separately, a
neighboring source may appear in more than one of them. Left alone, each
measurement fits that source independently, and the fits disagree.

The registry fixes that. It is one JSON file recording the fitted shape,
position and per-band brightness of shared sources, in sky units so an
entry works from any field.

Consume it:

\begin{cmd}\ttfamily\frenchspacing
sedphot measure -{}-name "M87" -{}-all -{}-out-dir M87 \textbackslash\\
\mbox{}~~~~-{}-registry registry.json
\end{cmd}

Consume it and write this galaxy's results back into it:

\begin{cmd}\ttfamily\frenchspacing
sedphot measure -{}-name "M87" -{}-all -{}-out-dir M87 \textbackslash\\
\mbox{}~~~~-{}-registry registry.json -{}-registry-update
\end{cmd}

Two things to know before using it:

\begin{itemize}
  \item \textbf{Run galaxies one at a time.} Updates are last-writer-wins
        with no file locking. Two runs writing at once will corrupt the
        file.
  \item \textbf{Order matters while the file fills.} A galaxy measured
        after its neighbor was registered gets that neighbor's shape
        frozen in rather than fitted fresh. Once every galaxy in a field
        has been measured once, a second pass over the same file is
        uniform.
\end{itemize}

Use it when you are measuring several galaxies in one field and want their
shared neighbors treated identically. Skip it for a single galaxy; it buys
nothing there.

\begin{notebox}
\emph{Every number this package writes is accompanied by a
\textup{\texttt{.provenance.json}} file recording the settings, the code
revision, and a checksum of the product. When a result needs defending,
that file is the record.}
\end{notebox}

\end{document}
